%% ma: L11-B10-0005
%% lop: 11
%% chuong: 4
%% bai: 10
%% dang: 11.10.6
%% loai: TLN
%% muc: VD
%% dap_an: "2"
%% nguon: "de-cuong-giua-ky-1-toan-11-nam-2026-2027-truong-thpt-hoang-van-thu-ha-noi-2-8.pdf (M08, Đề cương giữa kì 1 Toán 11 THPT Hoàng Văn Thụ, Hà Nội), Phần 3 Câu 8"
%% kien_thuc: "Bài 10 (giao điểm của đường thẳng và mặt phẳng), định lí Thalès (lớp 8)"
%% trang_thai: cho-duyet
%% ly_do: "Dạng mới đề xuất 11.10.6"
\begin{ex}
Cho hình chóp $S.ABCD$ có đáy $ABCD$ là hình bình hành. Gọi $M$, $P$ lần lượt là trung điểm của các cạnh $SA$ và $SC$. Điểm $N$ thuộc cạnh $SB$ sao cho $\dfrac{SN}{SB}=\dfrac23$. Gọi $Q$ là giao điểm của cạnh $SD$ và mặt phẳng $(MNP)$. Biết tỉ số $\dfrac{SQ}{SD}=\dfrac{a}{5}$. Tìm $a$.
\shortans{2}
\loigiai{Gọi $O=AC\cap BD$, $I=MP\cap SO$. Vì $MP$ là đường trung bình của $\triangle SAC$ nên $I$ là trung điểm $SO$; $Q=NI\cap SD$.

Trong $\triangle SBD$, qua $B$, $D$ kẻ các đường thẳng song song với $NQ$ cắt $SO$ tại $B'$, $D'$. Do $O$ là trung điểm $BD$ nên $O$ là trung điểm $B'D'$, suy ra $SB'+SD'=2SO=4SI$. Theo định lí Thalès: $\dfrac{SB}{SN}=\dfrac{SB'}{SI}$, $\dfrac{SD}{SQ}=\dfrac{SD'}{SI}$, nên $\dfrac{SB}{SN}+\dfrac{SD}{SQ}=4$.

Do đó $\dfrac32+\dfrac{SD}{SQ}=4\Rightarrow\dfrac{SQ}{SD}=\dfrac25$, $a=2$.}
\end{ex}
